\documentclass[11pt,a4paper]{article}

\usepackage[utf8]{inputenc}
\usepackage[T1]{fontenc}
\usepackage{lmodern}
\usepackage[margin=1.0in,headheight=14.5pt]{geometry}
\usepackage{microtype}
\usepackage{setspace}
\setstretch{1.14}
\emergencystretch=2.0em
\hbadness=10000

\usepackage{amsmath,amssymb,amsfonts,amsthm}
\usepackage{mathtools}
\usepackage{bm}
\usepackage{graphicx}
\usepackage{multirow}
\usepackage{float}
\usepackage{booktabs}
\usepackage{array}
\usepackage{tabularx}
\usepackage{enumitem}
\usepackage{listings}
\usepackage{upquote}
\usepackage{xcolor}
\usepackage{tcolorbox}
\usepackage{tikz}
\usepackage{fancyhdr}
\usepackage[unicode=true]{hyperref}

\newtheorem{proposition}{Proposition}

\hypersetup{
    colorlinks=true,
    linkcolor=blue!70!black,
    citecolor=blue!70!black,
    urlcolor=blue!70!black,
    pdftitle={Lumina Moon Standard Optimizer Evaluation Methodology (LSO-Eval v1.1.0)},
    pdfauthor={Silver Moon (cloverxion) - Lumina Moon (cloverx-id)}
}

\definecolor{codebg}{rgb}{0.96,0.96,0.97}
\definecolor{codekw}{rgb}{0.0,0.2,0.6}
\definecolor{codestr}{rgb}{0.65,0.15,0.15}

\lstset{
    backgroundcolor=\color{codebg},
    basicstyle=\ttfamily\footnotesize\linespread{0.94}\selectfont,
    keywordstyle=\color{codekw}\bfseries,
    stringstyle=\color{codestr},
    columns=fullflexible,
    keepspaces=true,
    upquote=true,
    breaklines=true,
    breakatwhitespace=true,
    showstringspaces=false,
    frame=single,
    rulecolor=\color{black!20},
    numbers=left,
    numberstyle=\tiny\color{black!50},
    numbersep=7pt,
    xleftmargin=1.8em,
    xrightmargin=0.2em,
    aboveskip=0.6em,
    belowskip=0.6em,
    tabsize=4
}

\tcbuselibrary{skins,breakable}
\newtcolorbox{ReferenceCase}[1]{
    colback=blue!2!white,
    colframe=blue!40!black,
    fonttitle=\bfseries\small,
    title={Reference Case Example: #1},
    boxrule=0.8pt,
    arc=2pt,
    breakable,
    before upper={\sloppy\emergencystretch=3.0em\hbadness=10000}
}

\AddToHook{shipout/background}{%
  \IfFileExists{luminamoon.png}{%
    \begin{tikzpicture}[remember picture, overlay]
      \node[opacity=0.08] at (current page.center) {%
        \includegraphics[width=0.8\textwidth]{luminamoon.png}%
      };
    \end{tikzpicture}%
  }{}%
}

\pagestyle{fancy}
\fancyhf{}

\renewcommand{\headrulewidth}{0.4pt}
\renewcommand{\footrulewidth}{0.4pt}

\fancyhead[L]{\small\textbf{LSO-Eval} \textbar\ Engineering Standard Specification}
\fancyhead[R]{\small v1.1.0 \textbar\ Apache-2.0}

\fancyfoot[L]{\footnotesize Lumina Moon (cloverx-id) \textbar\ Open Technical Specification}
\fancyfoot[R]{\footnotesize Page \thepage}

\begin{document}

\title{\textbf{Lumina Moon Standard Optimizer Evaluation Methodology (LSO-Eval)}\\[0.3em]
\large An Engineering Standard and Operational Procedure for Neural Network Optimization Algorithms}
\author{\textbf{Silver Moon} (\texttt{cloverxion})\\[0.2em]
\normalsize Lumina Moon (\texttt{cloverx-id})}
\date{October 2026 \\ \small Engineering Standard Specification --- Version 1.1.0}
\maketitle

\begin{abstract}
\noindent This manual defines the official \textbf{Lumina Moon Standard Optimizer Evaluation Methodology (LSO-Eval v1.1.0)}, an algorithm-agnostic verification and benchmarking standard for qualifying neural network optimizers and hardware acceleration backends. It establishes a repeatable, four-stage testing protocol to evaluate first-order, state-compressed, and quasi-second-order optimization algorithms alongside hardware execution backends (Rust AVX-512/AVX2 SIMD, Vulkan 1.3/1.4 Wave32 UMA, Triton AOT/JIT, and zero-copy NumPy/PyTorch frontends) before production deployment. LSO-Eval replaces ad-hoc trial-and-error with a structured four-stage filtering pipeline: (I) Formal Logic, Rational Pad\'{e} Approximation, and Bitwise Precision Verification, (II) Synthetic Stress-Test Surfaces, Cross-Backend Numerical Parity, and State Serialization Integrity, (III) Short-Run Proxy Model Screening, and (IV) Large-Scale Pre-Training Qualification with Cryptographic Zero-Leakage Validation and Resource-Tiered Pareto Acceptance. Reference case studies throughout this specification illustrate the standard using the \textbf{LuminaV} optimizer family (\texttt{v1.5.3}--\texttt{v1.5.5}).
\end{abstract}

\vspace{0.4em}
\hrule
\vspace{1.2em}

\section{Scope, Principles, and the Staged Evaluation Pipeline}

\subsection{Objective and Operational Scope}
The LSO-Eval standard governs the qualification of all first-order, state-compressed, and quasi-second-order optimization algorithms and hardware acceleration kernels developed or deployed within the Lumina Moon research and production infrastructure. Because LSO-Eval is an independent evaluation standard, its version lifecycle (\texttt{v1.1.0}) is decoupled from the version numbers of specific candidate algorithms. The objective is four-fold:
\begin{enumerate}[noitemsep,topsep=2pt]
    \item \textbf{Prevent Compute Wastage:} Filter out mathematically flawed, numerically unstable, or overly sensitive candidates at near-zero compute cost before allocating GPU clusters.
    \item \textbf{Separate Implementation Bugs from Theory:} Isolate algebraic flaws, rational function approximation errors, low-precision stochastic rounding bugs, and hardware kernel overhead into distinct test phases.
    \item \textbf{Guarantee Cross-Backend Consistency:} Enforce strict numerical parity across Scalar Reference Implementations, x86\_64 AVX-512/AVX2 SIMD, Vulkan 1.3/1.4 UMA Wave32 compute shaders, Triton AOT/JIT GPU kernels, and NumPy/PyTorch bindings (\texttt{PyO3 abi3/abi3t}).
    \item \textbf{Enforce Fair Multi-Objective Non-Regression:} Evaluate candidate optimizers across both validation quality and isolated hardware footprint (separating pure optimizer state memory from fixed model and activation baselines) while preventing regressions across releases.
\end{enumerate}

\subsection{Staged Filtering Protocol}
Candidate algorithms must pass through four sequential stages. Failure at any stage halts the evaluation immediately; engineering teams must not proceed to large-scale training while lower-stage checks remain unverified.

\begin{table}[!htbp]
\centering
\caption{The LSO-Eval Staged Evaluation Pipeline (v1.1.0)}
\label{tab:lso_funnel}
\vspace{0.4em}
\renewcommand{\arraystretch}{1.25}
\begin{tabularx}{\textwidth}{@{}l >{\raggedright\arraybackslash}p{3.2cm} >{\raggedright\arraybackslash}X c@{}}
\toprule
\textbf{Stage} & \textbf{Designation} & \textbf{Core Verification Objective} & \textbf{Cost Ceiling} \\
\midrule
\textbf{Stage I} & Math \& Bitwise Check & Zero division-by-zero, unit velocity bounds, subnormal floor safety, $[13/6]$ Pad\'{e} $\tanh$ accuracy, zero stochastic rounding bias. & $\mathcal{O}(\text{ms})$ \\
\textbf{Stage II} & Synthetic \& Parity & Ill-conditioned traversal ($\kappa \ge 1000$), saddle escape, $\lambda^*$ balance, 1-Pass EMA fusion, cross-backend numerical parity, 6-slot scratch \& checkpoint integrity. & $\mathcal{O}(\text{sec})$ \\
\textbf{Stage III} & Short-Run Proxy & Matrix update-to-weight ratio stability ($R_{t,l} \in [10^{-7}, 0.10]$), spike recovery ($\text{F}_2$), low-precision stall gating ($\text{F}_3 < 0.25$). & $\mathcal{O}(\text{min})$ \\
\textbf{Stage IV} & Large-Scale Run & Zero-leakage WAL validation, low-precision BF16/FP16 stability, isolated optimizer VRAM ($\text{VRAM}_{\text{opt}}$) vs.\ peak VRAM, kernel throughput, Two-Tiered Pareto held-out perplexity. & $\mathcal{O}(\text{hours})$ \\
\midrule
\textbf{D-Eval} & Version Regression & 4-Gate comparative audit evaluating release candidate $v_{k}$ against baseline predecessor $v_{k-1}$. & Cross-Stage \\
\bottomrule
\end{tabularx}
\end{table}

\section{Stage I: Formal Logic \& Symbolic Verification Protocol}

Stage I requires formal verification of algebraic safety properties and bitwise precision bounds before running any tensor training loop. The candidate optimizer is expressed in symbolic real arithmetic and checked using automated theorem provers and statistical equivalence tests.

\subsection{Standard SMT Safety Checklist (Z3 Solver)}
The engineering team formulates the optimizer update velocity $u_t(\theta, g, m, v)$ and checks seven core safety categories (implemented via twelve negated sub-queries in the reference harness) using the Z3 Theorem Prover~\cite{demoura2008z3}. Each property is verified by proving that its \textbf{negation} is \texttt{UNSAT} (unsatisfiable):

\begin{enumerate}[label=\textbf{Q\arabic*.}, itemsep=3pt]
    \item \textbf{Denominator Non-Singularity (Division-by-Zero):} For any gradient $g \in \mathbb{R}$, momentum $m \in \mathbb{R}$, and instantaneous defect residual $d_t = g - m$, prove that the envelope denominator $\mathcal{E}$ (or variance divisor $\sigma$) strictly satisfies $\mathcal{E} > 0$~\cite{maligranda2008triangle}. Negated query: $\mathcal{E} \le 0$.
    \item \textbf{Velocity Regularity, Trust-Region Radii \& Profile Boundedness:} Prove that the update velocity and trust-region governors satisfy their design profile:
    \begin{itemize}[noitemsep,topsep=2pt]
        \item \textit{Profile A (Coordinate-Bounded / Normalized, e.g., LuminaV, Lion~\cite{chen2023lion}):} Strictly bounded magnitude $\lvert u_t \rvert < 1.0$, positive coordinate trust radius $B_i^2 = r_{\text{bound}}^2(\theta_i^2 + 10^{-4}) > 0$, and radial trust lower bound $R = r_{\text{bound}}\max(\lVert\theta\rVert_2, 1.0) \ge r_{\text{bound}}$~\cite{brock2021agc}. Negated queries: $\lvert u_t \rvert \ge 1.0$, $B_i^2 \le 0$, or $R < r_{\text{bound}}$.
        \item \textit{Profile B (Variance-Adaptive, e.g., AdamW~\cite{loshchilov2019adamw}, AdaBelief~\cite{zhuang2020adabelief}, ADOPT~\cite{taniguchi2024adopt}):} Clamped variance divisor strictly positive $\sigma_{\text{clamped}} = \max(c_1\sqrt{v_t} + c_2, 10^{-12}) > 0$ and non-divergent quotient as $g \to 0$. Negated query: $\sigma \le 0 \lor \sigma_{\text{clamped}} \le 0$.
        \item \textit{Profile C (Matrix-Preconditioned / Subspace-Projected, e.g., Shampoo~\cite{gupta2018shampoo}, GaLore~\cite{zhao2024galore}):} Spectral/operator norm boundedness $\lVert U_t \rVert_2 \le C$ and non-singular preconditioning $\sigma_{\min}(P) > 0$.
    \end{itemize}
    \item \textbf{Radial Scaling \& Cautious Rescaling Bounds:} For any rational radial scaling factor $S(r)$ (such as the Pad\'{e} bound $(15 + r^2)/(15 + 6r^2)$ for $r > 0$ and $1.0$ at $r = 0$), prove that $0 < S(r) \le 1.0$ for all $r \ge 0$. For cautious alignment masking~\cite{liang2026cautious} with density floor $c_{\min} \in (0, 1]$, prove that the rescaling multiplier $\gamma_{\text{cautious}} = 1 / \mathrm{clamp}(\bar{m}, c_{\min}, 1.0)$ satisfies $\gamma_{\text{cautious}} \in [1.0, 1/c_{\min}]$.
    \item \textbf{Half-Precision Safe-Zone \& Subnormal Numerical Floor:} For half-precision execution (FP16 with in-container Bitwise-BF16 state promotion and unified $\epsilon_{\text{eff}} = \epsilon = 10^{-8} > 0$), prove that parameter update safety bounds remain strictly above the IEEE 754 FP16 subnormal threshold ($\text{floor}_{\text{sub}} = 6.1035 \times 10^{-5}$): $b_{\text{safe}} = \max(0.01\lvert\theta_0\rvert, 1.5 \times 10^{-4}) \ge 1.5 \times 10^{-4} > \text{floor}_{\text{sub}}$, while FP16 variance clamping $v_{\text{fp16}} = \min(v_{\text{raw}}, 65{,}000.0)$ remains strictly below the FP16 overflow ceiling ($65{,}504.0$).
    \item \textbf{Descent Direction Guarantee:} Prove that under zero initial momentum ($m_0 = 0$) or cautious alignment masking~\cite{liang2026cautious} ($m_{\text{mask}} = \mathbb{I}(u_t \cdot g_t > 0)$ with $\gamma_{\text{cautious}} \ge 1.0$), the effective update preserves non-negative alignment with the gradient: $u_0 \cdot g_t \ge 0$ and $\delta_{\text{cautious}} \cdot g_t \ge 0$.
    \item \textbf{Weight Decay Contraction:} For base learning rate $\eta \in (0, 0.05]$ and decoupled weight decay~\cite{loshchilov2019adamw} $\lambda_{\text{wd}} \in [0, 0.2]$, prove that $(1 - \eta \lambda_{\text{wd}}) \in (0, 1]$. Negated query: $(1 - \eta \lambda_{\text{wd}}) \le 0 \lor (1 - \eta \lambda_{\text{wd}}) > 1$.
    \item \textbf{Lookahead Bias Correction Positivity:} For all integer steps $t \ge 1$ and $\beta_1, \beta_2 \in (0, 1)$ with $0 < \beta_1^{t+1} \le \beta_1^2$ and $0 < \beta_2^t \le \beta_2$, prove that the Nesterov lookahead bias correction~\cite{zhang2019lookahead} $\mathrm{bc}_{1,\text{nes}} = 1 - \beta_1^{t+1}$ and clamped second-moment correction $\mathrm{bc}_2 = \max(1 - \beta_2^t, 10^{-15})$ remain strictly positive. Negated query: $\mathrm{bc}_{1,\text{nes}} \le 0 \lor \mathrm{bc}_2 \le 0$.
\end{enumerate}

\subsection{Standard Z3 SMT Verification Harness}
Listing~\ref{lst:z3_harness} provides the executable 12-query Stage I test script (\texttt{eval/modules/smt.py}) verifying Single-Buffer TD-CLM and Dual-Buffer Convex Lookahead Momentum (CLM) invariants.

\begin{lstlisting}[language=Python, caption={Standard LSO-Eval Stage I.1: 12-Query Z3 SMT Verification Harness (v1.1.0)}, label={lst:z3_harness}]
from z3 import If, Or, Real, Reals, Solver, unsat

def run_z3_smt_audit() -> bool:
    solver = Solver()
    m, g, w1, w2, lam, eps = Reals("m g w1 w2 lam eps")
    base_conds = [w1 > 0, w2 > 0, (w1 + w2) == 1, lam >= 0, eps > 0]
    abs_m, abs_g = If(m >= 0, m, -m), If(g >= 0, g, -g)
    diff = g - m
    abs_diff = If(diff >= 0, diff, -diff)
    clm_m = w1 * m + w2 * g
    clm_env = w1 * abs_m + w2 * abs_g + lam * abs_diff + eps
    abs_clm_m = If(clm_m >= 0, clm_m, -clm_m)

    # 1-2. TD-CLM Non-Singularity (Q1) & Unit Metric Bound |u| < 1.0 (Q2)
    solver.push(); solver.add(*base_conds, clm_env <= 0)
    res_tdclm_div0 = solver.check(); solver.pop()
    solver.push(); solver.add(*base_conds, abs_clm_m >= clm_env)
    res_tdclm_bound = solver.check(); solver.pop()

    # 3-4. Buffer-2 Sigma Positivity (Q1/Q2) & Cautious Scale [1, 1/c_min] (Q3)
    sqrt_v, c1, c2, m_bar, c_min = Reals("sqrt_v c1 c2 m_bar c_min")
    sigma = sqrt_v * c1 + c2
    sigma_clamped = If(sigma >= 1e-12, sigma, 1e-12)
    solver.push(); solver.add(sqrt_v >= 0, c1 > 0, c2 > 0, Or(sigma <= 0, sigma_clamped <= 0))
    res_sigma = solver.check(); solver.pop()
    solver.push(); solver.add(c_min == 0.5, m_bar >= c_min, m_bar <= 1.0)
    solver.add(Or((1.0 / m_bar) < 1.0, (1.0 / m_bar) > (1.0 / c_min)))
    res_cautious = solver.check(); solver.pop()

    # 5-6. Coordinate Governor B > 0 & Radial Lower Bound R >= B_ratio (Q2)
    p, B_ratio, norm_p = Reals("p B_ratio norm_p")
    solver.push(); solver.add(B_ratio > 0, (B_ratio * B_ratio) * (p * p + 1e-4) <= 0)
    res_coord_b = solver.check(); solver.pop()
    solver.push(); solver.add(B_ratio > 0, norm_p >= 0)
    R = If(norm_p >= 1.0, B_ratio * norm_p, B_ratio * 1.0)
    solver.add(R < B_ratio); res_radial_r = solver.check(); solver.pop()

    # 7-8. FP16 Safe-Zone Floor >= 1.5e-4 & Subnormal/Variance Safety (Q4)
    p_fp16, v_raw = Reals("p_fp16 v_raw")
    p_scaled = 0.01 * If(p_fp16 >= 0, p_fp16, -p_fp16)
    b_safe = If(p_scaled >= 1.5e-4, p_scaled, 1.5e-4)
    v_fp16 = If(v_raw <= 65000.0, v_raw, 65000.0)
    solver.push(); solver.add(b_safe < 1.5e-4); res_fp16 = solver.check(); solver.pop()
    solver.push(); solver.add(Or(b_safe <= 6.1035e-5, v_fp16 > 65504.0))
    res_subnormal = solver.check(); solver.pop()

    # 9-12. Radial Scale (Q3), Bias Corr (Q7), Weight Decay (Q6), Descent (Q5)
    r, lr, wd = Reals("r lr wd")
    solver.push(); solver.add(r >= 0)
    s_radial = If(r == 0, 1.0, (15.0 + r * r) / (15.0 + 6.0 * r * r))
    solver.add(Or(s_radial <= 0.0, s_radial > 1.0)); res_radial_scale = solver.check(); solver.pop()

    beta1, beta2, b1_tp1, b2_t = Reals("beta1 beta2 b1_tp1 b2_t")
    solver.push(); solver.add(beta1 > 0, beta1 < 1, beta2 > 0, beta2 < 1,
               b1_tp1 > 0, b1_tp1 <= beta1 * beta1, b2_t > 0, b2_t <= beta2)
    bc2_clamped = If(1.0 - b2_t >= 1e-15, 1.0 - b2_t, 1e-15)
    solver.add(Or((1.0 - b1_tp1) <= 0, bc2_clamped <= 0)); res_bias_corr = solver.check(); solver.pop()

    solver.push(); solver.add(lr > 0, lr <= 0.05, wd >= 0, wd <= 0.2)
    solver.add(Or((1.0 - lr * wd) <= 0, (1.0 - lr * wd) > 1)); res_wd = solver.check(); solver.pop()

    g_val, u_val, w2_val, lam_val, eps_val, gamma_c = Reals("g_val u_val w2_val lam_val eps_val gamma_c")
    solver.push(); solver.add(w2_val > 0, lam_val >= 0, eps_val > 0, gamma_c >= 1.0)
    abs_g_val = If(g_val >= 0, g_val, -g_val)
    u_zero_m = (w2_val * g_val) / (w2_val * abs_g_val + lam_val * abs_g_val + eps_val)
    delta_cautious = u_val * If((u_val * g_val) > 0, 1.0, 0.0) * gamma_c
    solver.add(Or(u_zero_m * g_val < 0, delta_cautious * g_val < 0)); res_descent = solver.check(); solver.pop()

    return all(s == unsat for s in [
        res_tdclm_div0, res_tdclm_bound, res_sigma, res_cautious, res_coord_b, res_radial_r,
        res_fp16, res_subnormal, res_radial_scale, res_bias_corr, res_wd, res_descent
    ])
\end{lstlisting}

\subsection{Symbolic Calculus, Rational Pad\'{e} Audit \& Bitwise Precision}
After passing SMT verification, the optimizer undergoes four analytical and bitwise precision checks (implemented via nine automated tests in \texttt{eval/modules/math\_verification.py}):
\begin{enumerate}[label=\textbf{S\arabic*.}, itemsep=3pt]
    \item \textbf{Large-Gradient Limit \& Equilibrium Check:} Using SymPy~\cite{meurer2017sympy}, evaluate the symbolic limit of velocity $u_t$ as gradient magnitude approaches infinity:
    \begin{equation}
        u_{\infty} = \lim_{|g| \to \infty} u_t(m, g).
    \end{equation}
    \textbf{Acceptance Rule:} For Profile A (Coordinate-Bounded), the limit must remain strictly below unit magnitude ($\lvert u_\infty \rvert < 1.0$) and match the closed-form spike--reversal equilibrium $\lambda^*$. For Profile B and C, the limit must remain finite without numerical overflow ($\lim_{\lvert g \rvert \to \infty} \lVert u_t \rVert < \infty$).
    \item \textbf{Piecewise Continuity ($C^0/C^1$ Transition) \& Derivative Bounds:} When radial bounding switches from unity scaling ($S(r) = 1.0$ for $r < r_{\text{thresh}}$) to $\tanh(r)/r$ at threshold $r_{\text{thresh}} = 10^{-4}$, compute the transition step:
    \begin{equation}
        \Delta_{\text{jump}} = \left| 1.0 - \frac{\tanh(r_{\text{thresh}})}{r_{\text{thresh}}} \right| \approx \frac{r_{\text{thresh}}^2}{3}.
    \end{equation}
    Similarly, for coordinate bounding $g(x) = B\tanh(x/B)$, verify that $g'(0) = 1.0$ and $\lim_{x\to\infty} g'(x) = 0.0$.
    \textbf{Acceptance Rule:} $\Delta_{\text{jump}}$ must be strictly smaller than single-precision machine epsilon ($\epsilon_{\text{FP32}} \approx 1.19 \times 10^{-7}$). At $r_{\text{thresh}} = 10^{-4}$, $\Delta_{\text{jump}} \approx 3.33 \times 10^{-9} \ll \epsilon_{\text{FP32}}$.
    \item \textbf{Rational Pad\'{e} $[13/6]$ Approximation Accuracy:} For hardware backends without native single-cycle $\tanh$ instructions (x86\_64 AVX-512/AVX2 and Vulkan SPIR-V), $\tanh(x)$ is evaluated via a pole-free clamped $[13/6]$ rational Pad\'{e} approximant~\cite{molina2020pade} on $x \in [-7.9988117, 7.9988117]$:
    \begin{equation}
        \tanh_{\text{Pad\'{e}}}(x) = \frac{x \sum_{k=0}^{6} \alpha_{2k+1} x^{2k}}{\sum_{j=0}^{3} \beta_{2j} x^{2j}}.
    \end{equation}
    \textbf{Acceptance Rule:} Maximum absolute error across $\mathbb{R}$ against IEEE 754 \texttt{std::f32::tanh} must satisfy $\|\tanh_{\text{Pad\'{e}}}(x) - \tanh(x)\|_\infty < 1.5 \times 10^{-6}$, while preserving exact odd symmetry $\tanh_{\text{Pad\'{e}}}(-x) = -\tanh_{\text{Pad\'{e}}}(x)$ and IEEE NaN propagation.
    \item \textbf{Stochastic Rounding Unbiasedness \& MurmurHash3 Determinism:} Low-precision quantization (\texttt{bfloat16} and \texttt{float16})~\cite{gupta2015limited, zamirai2021bfloat16} uses the stateless MurmurHash3 32-bit finalizer seeded deterministically per parameter and step:
    \begin{align}
        \text{seed}_m &= (\text{seed}_0 + t \cdot \mathtt{0x9E3779B9} + \text{idx}_{\text{param}} \cdot 10007) \ \& \ \mathtt{0x7FFFFFFF}, \\
        \text{seed}_p &= (\text{seed}_m + \mathtt{0x55555555}) \ \& \ \mathtt{0x7FFFFFFF}.
    \end{align}
    For FP16, the bitmask width adapts to subnormal exponents ($\mathrm{exp}_{\text{f32}} < 113$) via $\mathrm{shift}_{\text{total}} = 13 + \mathrm{clamp}(113 - \mathrm{exp}_{\text{f32}}, 0, 11)$. Generate $N = 200{,}000$ samples across both normal ($[10^{-3}, 10^{-1}]$) and FP16 subnormal ($[10^{-5}, 6.1 \times 10^{-5}]$) ranges. Evaluate unbiasedness via a two-sided Student's $t$-test ($H_0: E[\mathcal{Q}_{\text{SR}}(x)] = x$) and Two One-Sided Tests (TOST) within tolerance $\pm \delta_{\text{tol}} = 10^{-6}$.
    \textbf{Acceptance Rule:} Maximum mean bias must satisfy $\lvert \mu_{\mathcal{Q}} - x \rvert < 10^{-6}$ ($p_{\text{TOST}} < 0.01$, or $t$-test $p \ge 0.01$) across both normal and subnormal ranges.
\end{enumerate}

\begin{ReferenceCase}{LuminaV (v1.5.3--v1.5.5) TD-CLM \& Dual-Buffer CLM Verification in Stage I}
In LuminaV (\texttt{v1.5.3} uncentered mode through \texttt{v1.5.5} canonical specification), the first moment updates via $m_t = \beta_1 m_{t-1} + (1 - \beta_1) g_t$. Single-Buffer TD-CLM (\texttt{buffer=1}) and Dual-Buffer Convex Lookahead Momentum (\texttt{buffer=2}) compute raw update velocity via convex lookahead interpolation~\cite{zhang2019lookahead}, triangle-defect norm envelopes~\cite{maligranda2008triangle}, and decorrelated belief variance~\cite{zhuang2020adabelief, taniguchi2024adopt}:
\begin{align*}
    u_t^{(\text{B1})} &= \frac{w_1 m_{t-1} + w_2 g_t}{w_1 |m_{t-1}| + w_2 |g_t| + \lambda |d_t^{(\text{B1})}| + \epsilon_{\text{eff}}}, \quad w_1 = \beta_1^2, \quad w_2 = 1 - \beta_1^2, \\
    u_t^{(\text{B2})} &= \tanh\left(\frac{\beta_1 m_t + (1 - \beta_1) g_t}{\max(c_1 \sqrt{\max(v_t, 0)} + c_2, 10^{-12})}\right), \quad v_t = \beta_2 v_{t-1} + (1 - \beta_2)\left(d_t^{(\text{B2})}\right)^2,
\end{align*}
where the coordinate-wise innovation defects and lookahead bias-correction scalars are defined as:
\begin{align*}
    d_t^{(\text{B1})} &= g_t - m_{t-1}, \quad d_t^{(\text{B2})} = g_t - m_t, \\
    c_1 &= \frac{(1 - \beta_1^{t+1})\tau}{\sqrt{\max(1 - \beta_2^t, 10^{-15})}}, \quad c_2 = \epsilon_{\text{eff}} (1 - \beta_1^{t+1})\tau.
\end{align*}
Note the algebraic identity $\beta_1 m_t + (1 - \beta_1)g_t \equiv \beta_1^2 m_{t-1} + (1 - \beta_1^2)g_t = w_1 m_{t-1} + w_2 g_t$, unifying the Nesterov lookahead numerator across both buffer modes.
\begin{enumerate}[label=\textbf{\arabic*.}, noitemsep, topsep=2pt]
    \item \textbf{12-Query Z3 SMT Result:} By the triangle inequality~\cite{maligranda2008triangle}, $\lvert w_1 m_{t-1} + w_2 g_t \rvert \le w_1 \lvert m_{t-1} \rvert + w_2 \lvert g_t \rvert$. Because $\lambda \lvert d_t^{(\text{B1})} \rvert \ge 0$ and $\epsilon_{\text{eff}} = 10^{-8} > 0$, the Buffer 1 denominator strictly exceeds the numerator, while Buffer 2 satisfies $\lvert \tanh(\cdot) \rvert < 1.0$ with denominator floor $10^{-12} > 0$. All 12 negated queries in Listing~\ref{lst:z3_harness} returned \texttt{PASSED (UNSAT)}.
    \item \textbf{Large-Gradient Ceiling \& Bias Correction:} SymPy limit analysis on Buffer 1 yields $u_{\infty} = \frac{w_2}{w_2 + \lambda} = \frac{0.19}{0.49} \approx 0.3878 < 1.0$ and equilibrium $\lambda^* = \frac{2(0.19)^2}{0.24} \approx 0.3008$. Buffer 2 bias correction $c_1$ evaluates to $4.8067$ at $t = 1$ and converges monotonically to $\tau = 0.8000$ as $t \to \infty$.
    \item \textbf{Pad\'{e} $[13/6]$ \& Stochastic Rounding Precision:} Across $500{,}000$ evaluation points on $[-8, 8]$, $[13/6]$ Pad\'{e} $\tanh$ achieved maximum absolute error $\mathbf{9.48 \times 10^{-8}} \ll 1.5 \times 10^{-6}$ with exact odd symmetry. Monte Carlo evaluation over $200{,}000$ samples of MurmurHash3 bitwise rounding across normal ($\pm 0.12345678$) and FP16 subnormal ($\pm 3.0 \times 10^{-5}$) values yielded maximum absolute bias $\mathbf{3.53 \times 10^{-7}} < 10^{-6}$ ($p_{\text{TOST}} = \mathbf{0.00} < 0.01$).
\end{enumerate}
\end{ReferenceCase}

\section{Stage II: Synthetic Test Surfaces, Cross-Backend Parity \& State Serialization}

Stage II connects mathematical verification (Stage I) to neural network training. Its goal is to test optimizer trajectories on challenging synthetic surfaces, calibrate default hyperparameters, verify numerical consistency across hardware backends, and check state serialization integrity.

\subsection{Synthetic Benchmark Functions}
The engineering team evaluates the candidate optimizer against reference baselines (standard AdamW~\cite{loshchilov2019adamw} in both FP32 Master and Native BF16 precision, plus the previous algorithm baseline) across five standard test functions~\cite{jamil2013benchmark}:

\begin{enumerate}[label=\textbf{B\arabic*.}, itemsep=3pt]
    \item \textbf{Ill-Conditioned Quadratic Bowl ($\kappa \ge 1000$):}
    \begin{equation}
        f(w_0, w_1) = \frac{1}{2}\left(w_0^2 + 1000 w_1^2\right), \quad w^{(0)} = [10.0, 1.0]^\top.
    \end{equation}
    \textbf{Evaluation Focus:} Tests whether the triangle-defect envelope~\cite{maligranda2008triangle} or belief variance~\cite{zhuang2020adabelief} dampens steep-wall oscillations along $w_1$ without slowing progress along the flat direction $w_0$.
    \item \textbf{Non-Convex Rosenbrock Curved Valley~\cite{jamil2013benchmark, rosenbrock1960}:}
    \begin{equation}
        f(w_0, w_1) = (1 - w_0)^2 + 100(w_1 - w_0^2)^2, \quad w^{(0)} = [-1.5, 2.0]^\top.
    \end{equation}
    \textbf{Evaluation Focus:} Tests the optimizer's ability to follow curved non-linear valleys where the gradient direction shifts continuously.
    \item \textbf{Stochastic Gradient Noise Resilience:}
    Evaluates the quadratic bowl under additive Gaussian gradient noise $\tilde{g}_t = \nabla f(w_t) + \mathcal{N}(0, \sigma^2 I)$ with $\sigma = 1.0$.
    \textbf{Evaluation Focus:} Measures steady-state error and noise amplification near the optimum.
    \item \textbf{Multimodal Rastrigin Local-Minima Escape~\cite{rastrigin1974}:}
    \begin{equation}
        f(w_0, w_1) = 20 + \sum_{i=0}^1 \left(w_i^2 - 10\cos(2\pi w_i)\right), \quad w^{(0)} = [2.5, 2.5]^\top.
    \end{equation}
    \textbf{Evaluation Focus:} Evaluates whether lookahead momentum~\cite{zhang2019lookahead} carries the iterate out of shallow local minima toward the global minimum at $(0, 0)$.
    \item \textbf{High-Dimensional Saddle Point Escape ($N \ge 100$):}
    \begin{equation}
        f(w) = \frac{1}{2} \sum_{i=1}^N \lambda_i w_i^2 + \frac{1}{4} \sum_{i=1}^N w_i^4, \quad \lambda_i \in [-1.0, 1000.0].
    \end{equation}
    \textbf{Evaluation Focus:} Evaluates the optimizer's ability to escape high-dimensional saddle points (negative curvature eigenvalues) and verifies coordinate/radial bounding stability.
\end{enumerate}

\subsection{Standard 2D Simulation Harness}
Listing~\ref{lst:harness_2d} provides the deterministic Python harness for Stage II trajectory testing, implementing the TD-CLM update step with cautious rescaling~\cite{liang2026cautious} and coordinate bounding~\cite{brock2021agc}.

\begin{lstlisting}[language=Python, caption={Standard LSO-Eval Stage II: Deterministic Simulation Harness (v1.1.0)}, label={lst:harness_2d}]
import numpy as np

def run_2d_benchmark(loss_fn, grad_fn, step_fn, w_init, f_star=0.0,
                     lr=0.01, max_steps=2000, tol=1e-5, noise_std=0.0, seed=1337):
    rng = np.random.default_rng(seed)
    w = np.array(w_init, dtype=np.float64)
    state = {}
    current_loss = loss_fn(w)

    for t in range(1, max_steps + 1):
        g = grad_fn(w)
        if noise_std > 0.0:
            g = g + rng.normal(0.0, noise_std, size=g.shape)

        w, state = step_fn(w, g, state, lr, t)
        current_loss = loss_fn(w)
        curr_gap = abs(current_loss - f_star)

        if np.isnan(w).any() or np.isinf(w).any() or curr_gap > 1e8:
            return t, curr_gap, "DIVERGED"
        if curr_gap < tol:
            return t, curr_gap, "CONVERGED"

    return max_steps, abs(current_loss - f_star), "MAX_STEPS"

def luminav_tdclm_step_fn(w, g, state, lr, t,
                          beta1=0.9, lam=0.3, eps=1e-8, bound_ratio=0.05):
    if "m" not in state:
        state["m"] = np.zeros_like(w)

    m = state["m"]
    w1 = beta1 * beta1
    w2 = 1.0 - w1

    m_new = beta1 * m + (1.0 - beta1) * g
    clm_m = w1 * m + w2 * g
    diff = g - m
    clm_env = w1 * np.abs(m) + w2 * np.abs(g) + lam * np.abs(diff) + eps

    u = clm_m / clm_env
    mask = (u * g > 0.0).astype(np.float64)
    cautious_scale = 1.0 / np.clip(np.mean(mask), 0.5, 1.0)

    step_val = lr * (u * mask) * cautious_scale
    b_i = bound_ratio * np.sqrt(w * w + 1e-4)
    step_val = b_i * np.tanh(step_val / b_i)

    state["m"] = m_new
    w_next = w - step_val
    return w_next, state
\end{lstlisting}

\subsection{Hyperparameter Sensitivity \& Equilibrium Calibration}
Optimizers often fail in practice when their performance depends on a narrow hyperparameter window. LSO-Eval requires systematic sweeps across all hyperparameters:

\begin{enumerate}[label=\textbf{P\arabic*.}, itemsep=3pt]
    \item \textbf{One-at-a-Time Sweep Protocol:} Fix standard defaults ($\beta_1 = 0.9, \beta_2 = 0.999, \tau = 0.8, \epsilon = 10^{-8}$) and vary the target hyperparameter across at least 5 values spanning an order of magnitude.
    \item \textbf{Analytical Balance Point:} Whenever a parameter controls a trade-off (such as spike damping vs.\ sign-reversal responsiveness), derive the analytical balance point in closed form and enforce input validation bounds in code.
    \item \textbf{Parameter Stability Range:} The default value must sit within a stable range where convergence speed degrades by no more than $15\%$ under a $\pm 30\%$ perturbation ($\tau \in [0.56, 1.04]$ and $r_{\text{bound}} \in [0.035, 0.065]$).
\end{enumerate}

\subsection{Feature Combinations, 1-Pass EMA Parity \& Cross-Backend Numerical Parity}
Because production optimizer suites combine multiple algorithmic options and hardware execution backends, Stage II enforces three parity checks:

\begin{enumerate}[label=\textbf{A\arabic*.}, itemsep=3pt]
    \item \textbf{Feature Combination Stability:} Across all 12 factorial combinations of $\mathcal{C} = \mathcal{M}_{\text{buffer}} \times \mathcal{M}_{\text{bound}} \times \mathcal{M}_{\text{cautious}}$, enabling any subset of features must remain numerically finite and convergent. Coordinate bounding applies per-parameter adaptive clipping~\cite{brock2021agc} via soft hyperbolic saturation:
    \begin{equation}
        \Delta\theta_i = B_i \tanh\left(\frac{\eta \delta_i}{B_i}\right), \quad B_i = r_{\text{bound}} \sqrt{\theta_{0,i}^2 + 10^{-4}},
    \end{equation}
    while radial bounding applies unit-wise norm scaling $S(r) = \tanh(r)/r$ for $r = \frac{\eta \Vert{}\delta\Vert{}_2}{r_{\text{bound}}\max(\Vert{}\theta_0\Vert{}_2, 1.0)}$.
    \item \textbf{IO-Aware 1-Pass Fused Scratchpad EMA Parity:} Following IO-aware kernel fusion principles~\cite{dao2022flashattention, tillet2019triton}, in 1-Pass Fused mode (\texttt{fused\_\allowbreak single\_\allowbreak pass=True}, active for step $t > 1$) each parameter maintains a 6-element \texttt{float32} scratchpad vector (5 active slots + 1 reserved ABI alignment slot):
    \begin{equation}
        \mathbf{s} = \big[\underbrace{\Sigma m_{\text{mask}}}_{s_0},\; \underbrace{\Sigma u_{\text{eff}}^2}_{s_1},\; \underbrace{\Sigma \theta_0^2}_{s_2},\; \underbrace{\gamma_{\text{cautious}}}_{s_3},\; \underbrace{S_{\text{bound}}}_{s_4},\; \underbrace{0.0}_{s_5\text{ (reserved)}}\big] \in \mathbb{R}^6.
    \end{equation}
    The kernel reads $s_3$ and $s_4$ from step $t-1$ while accumulating $s_0, s_1, s_2$ via block-level register reductions and a single \texttt{atomic\_add} per thread block in one HBM memory pass, then updates $s_3, s_4$ post-step via EMA with mode-specific defaults (\texttt{ema\_\allowbreak decay="auto"} $\implies \alpha_{\text{ema}} = 0.95$ for \texttt{buffer=1} and $\alpha_{\text{ema}} = 0.80$ for \texttt{buffer=2}). The convergence step count of 1-Pass Fused EMA must not degrade by more than $+1.0\%$ relative to 2-Pass Instantaneous execution.
    \item \textbf{Cross-Backend Numerical Parity (\texttt{luminav-\allowbreak core}, \texttt{luminav-\allowbreak simd}, \texttt{luminav-\allowbreak vulkan}):}
    \begin{itemize}[noitemsep,topsep=2pt]
        \item \textit{Parallel Chunking \& Pairwise Tree Reduction:} Multithreaded Rayon execution partitions tensors into chunks of \texttt{ELEM\_\allowbreak CHUNK\_\allowbreak SIZE} $= 2{,}097{,}152$ elements using 16-lane pairwise tree reduction (\texttt{reduce\_\allowbreak sum\_16}). It must achieve exact bitwise parity (\texttt{atol} $= 0.0$) against single-threaded scalar execution.
        \item \textit{Bitwise PRNG \& Stochastic Rounding Parity:} AVX-512 and AVX2 vectorized PRNG and FP16/BF16 stochastic rounding must achieve exact \texttt{0.0} bit-level parity against the scalar reference implementation.
        \item \textit{Multi-Step Trajectory Parity ($N = 131{,}072$, 5 steps):} Across both TD-CLM and Dual-Buffer modes in FP32, BF16, and pure FP16 (with in-container Bitwise-BF16 state promotion), AVX-512, AVX2, and Vulkan 1.3/1.4 Wave32 UMA compute shaders must match the Scalar Reference within $\text{atol} \le 10^{-5}$ (FP32) and $\text{atol} \le 2 \times 10^{-2}$ (BF16/FP16).
        \item \textit{Frontend Wrapper Parity:} Standalone NumPy (\texttt{Standalone\allowbreak Optimizer} via zero-copy \texttt{\_\_array\_\allowbreak interface\_\_}) and PyTorch (\texttt{torch.\allowbreak optim.\allowbreak Optimizer}) must produce identical parameter trajectories on matching inputs.
    \end{itemize}
\end{enumerate}

\subsection{State Cleanup \& Checkpoint Serialization Audit}
To prevent silent resume bugs or memory leaks, the optimizer must pass the \textbf{LSO-Eval Serialization Audit} (Listing~\ref{lst:checkpoint_test}), verifying across both \texttt{float32} and \texttt{bfloat16} models and all master-weight modes (\texttt{False}, \texttt{"semi"}, \texttt{"full"}): (1) no temporary contiguous staging buffers remain in the state dictionary, (2) temporary reduction accumulators and reserved slots ($s_0, s_1, s_2, s_5$) are reset to zero while EMA scales ($s_3, s_4$) remain in \texttt{float32}, (3) \texttt{master\_\allowbreak param} retains \texttt{float32} precision, and (4) parameter updates after reloading match an uninterrupted run bit-for-bit.

\begin{lstlisting}[language=Python, caption={Standard LSO-Eval Stage II: Checkpoint Roundtrip \& State Cleanup Test (v1.1.0)}, label={lst:checkpoint_test}]
import copy, io, torch, torch.nn as nn

def verify_checkpoint_integrity(optimizer_cls):
    torch.manual_seed(42)
    forbidden_keys = (
        "p_contig", "exp_avg_contig", "exp_avg_sq_contig",
        "mp_contig", "_plan", "_contig_step"
    )
    for dtype in (torch.float32, torch.bfloat16):
        for buf in (1, 2):
            for mw in (False, "semi", "full"):
                model_cont = nn.Sequential(nn.Linear(10, 20), nn.ReLU(), nn.Linear(20, 2)).to(dtype=dtype)
                model_res = copy.deepcopy(model_cont)
                opt_cont = optimizer_cls(model_cont.parameters(), lr=1e-3, buffer=buf, bound=True, master_weights=mw)
                opt_res = optimizer_cls(model_res.parameters(), lr=1e-3, buffer=buf, bound=True, master_weights=mw)
                batches = [(torch.randn(4, 10, dtype=dtype), torch.tensor([0, 1, 0, 1])) for _ in range(6)]

                for i in range(5):
                    x, y = batches[i]
                    for m, opt in ((model_cont, opt_cont), (model_res, opt_res)):
                        loss = nn.CrossEntropyLoss()(m(x).float(), y)
                        opt.zero_grad(); loss.backward(); opt.step()

                buffer = io.BytesIO()
                torch.save(opt_res.state_dict(), buffer); buffer.seek(0)
                opt_loaded = optimizer_cls(model_res.parameters(), lr=1e-3, buffer=buf, bound=True, master_weights=mw)
                state_dict = torch.load(buffer, weights_only=False)
                opt_loaded.load_state_dict(state_dict)

                for p_state in state_dict["state"].values():
                    for k in forbidden_keys:
                        assert k not in p_state, f"Memory leak: '{k}' in state_dict!"
                    if "scratch" in p_state:
                        sc = p_state["scratch"]
                        assert float(sc[0]) == 0.0 and float(sc[1]) == 0.0
                        assert float(sc[2]) == 0.0 and float(sc[5]) == 0.0

                for p in model_res.parameters():
                    st = opt_loaded.state[p]
                    if "scratch" in st: assert st["scratch"].dtype == torch.float32
                    if "master_param" in st: assert st["master_param"].dtype == torch.float32

                x_last, y_last = batches[5]
                for m, opt in ((model_cont, opt_cont), (model_res, opt_loaded)):
                    loss = nn.CrossEntropyLoss()(m(x_last).float(), y_last)
                    opt.zero_grad(); loss.backward(); opt.step()

                for p_c, p_r in zip(model_cont.parameters(), model_res.parameters()):
                    torch.testing.assert_close(p_c, p_r, rtol=0.0, atol=0.0)
    return True
\end{lstlisting}

\begin{ReferenceCase}{LuminaV (v1.5.3--v1.5.5) Trajectory, Equilibrium \& Parity in Stage II}
\begin{proposition}[Symmetric Spike--Reversal Balance for $\lambda$]
\label{prop:lambda_equilibrium}
Let $w_1 = \beta_1^2$ and $w_2 = 1 - \beta_1^2$ with $\beta_1 \in (\frac{\sqrt{3}}{2}, 1)$. In TD-CLM~\cite{maligranda2008triangle}, the velocity ceiling under a large gradient spike ($\lvert g_t \rvert \to \infty$) is $u_{\text{spike}}(\lambda) = \frac{w_2}{w_2 + \lambda}$, and the instantaneous velocity ceiling under an exact gradient sign reversal ($g_t = -m_{t-1}$) is $u_{\text{rev}}(\lambda) = \frac{w_1 - w_2}{1 + 2\lambda}$. The unique coefficient $\lambda^*$ that equalizes spike damping and reversal damping ($u_{\text{spike}}(\lambda^*) = u_{\text{rev}}(\lambda^*)$) is given in closed form by:
\begin{equation}
    \lambda^* = \frac{2 w_2^2}{w_1 - 3w_2} = \frac{2(1 - \beta_1^2)^2}{4\beta_1^2 - 3}, \quad \forall \beta_1 \in \left(\frac{\sqrt{3}}{2}, 1\right).
\end{equation}
\end{proposition}
\begin{proof}
Setting $\frac{w_2}{w_2 + \lambda} = \frac{w_1 - w_2}{1 + 2\lambda}$ yields $w_2(1 + 2\lambda) = (w_1 - w_2)(w_2 + \lambda)$. Expanding and collecting terms in $\lambda$ gives $w_2(1 - w_1 + w_2) = (w_1 - 3w_2)\lambda$. Substituting $1 - w_1 = w_2$ produces $2w_2^2 = (w_1 - 3w_2)\lambda$, which solves directly to $\lambda^* = \frac{2(1-\beta_1^2)^2}{4\beta_1^2 - 3}$ for $4\beta_1^2 - 3 > 0 \iff \beta_1 > \frac{\sqrt{3}}{2} \approx 0.866$. Substituting standard $\beta_1 = 0.9$ yields $\lambda^* \approx 0.3008 \approx 0.30$.
\end{proof}

LuminaV (\texttt{v1.5.3}--\texttt{v1.5.5}) was verified across all Stage II checks:
\begin{itemize}[itemsep=2pt]
    \item \textbf{Synthetic Surfaces (\texttt{B1}--\texttt{B5}):}
    \begin{itemize}[noitemsep,topsep=1pt]
        \item \textit{B1 Quadratic ($\kappa=1000$, $\eta=0.02$, $1{,}000$ steps):} Unregularized B-CLM ($\lambda=0.0$) required $950$ steps, whereas \textbf{LuminaV Buffer 1 (TD-CLM)} converged in $\mathbf{575}$ \textbf{steps} (loss gap $2.67\times 10^{-6}$) and \textbf{LuminaV Buffer 2 (Dual)} in $\mathbf{571}$ \textbf{steps} ($9.31\times 10^{-6}$). Both \texttt{AdamW\_FP32\_Master} (gap $3.66\times 10^{-2}$) and \texttt{AdamW\_Native\_BF16} (stalled at gap $50.00$) failed to reach $10^{-5}$.
        \item \textit{B2 Rosenbrock ($\eta=0.002$, $2{,}000$ steps) \& B5 Saddle ($N=100$, $\eta=0.02$):} On Rosenbrock, Buffer 1 reached a loss gap of $7.46\times 10^{-4}$ and Buffer 2 reached $1.76$ vs.\ $3.90$ for \texttt{AdamW\_FP32\_Master} and $12.50$ for \texttt{AdamW\_Native\_BF16}. On $N=100$ Saddle Escape, \textbf{LuminaV Buffer 2} converged fastest in $\mathbf{149}$ \textbf{steps} ($5.25\times 10^{-6}$) vs.\ $193$ steps for \texttt{AdamW\_FP32\_Master} and $196$ steps for \texttt{AdamW\_Native\_BF16}.
        \item \textit{B3 Gaussian Noise ($\sigma=1.0$) \& B4 Rastrigin:} Under heavy gradient noise, Buffer 2 ($8.11\times 10^{-2}$) and Buffer 1 ($2.21\times 10^{-1}$) achieved $15\times$--$41\times$ lower steady-state loss than \texttt{AdamW\_FP32\_Master} ($3.36$) and \texttt{AdamW\_Native\_BF16} ($50.00$). On Rastrigin, Buffer 2 ($7.9597$) and Buffer 1 ($7.9690$) matched or beat \texttt{AdamW} ($7.9597$ / $7.9617$).
    \end{itemize}
    \item \textbf{Equilibrium $\lambda^* = 0.3008$ \& $\pm 30\%$ Basin Stability:} At $\lambda = 0.30$, spike ceiling ($0.3878$) and reversal ceiling ($0.3875$) equalized, converging in $575$ steps (while resonance $w_2 = \lambda = 0.19$ stalled). Perturbing $\tau \in [0.56, 1.04]$ (ref $556$ steps), coordinate bound $r_{\text{bound}} \in [0.035, 0.065]$ (ref $587$ steps), and radial bound $r_{\text{bound}} \in [0.035, 0.065]$ (ref $560$ steps) all passed within the $\le 15\%$ stability window.
    \item \textbf{1-Pass Fused Scratchpad EMA Parity:} On $\kappa=1000$, Buffer 1 1-Pass Fused EMA ($\alpha_{\text{ema}}=0.95$) converged in $\mathbf{346}$ \textbf{steps} vs.\ $381$ steps for 2-Pass Instantaneous ($\mathbf{-9.19\%}$ step reduction, loss $2.52\times 10^{-5}$), and Buffer 2 1-Pass Fused EMA ($\alpha_{\text{ema}}=0.80$) converged in $\mathbf{367}$ \textbf{steps} vs.\ $370$ steps ($\mathbf{-0.81\%}$ step reduction, loss $9.24\times 10^{-5}$).
    \item \textbf{Cross-Backend \& Serialization Parity:} All 12 factorial ablation combinations and all 24 checkpoint roundtrip configurations ($\{\text{FP32}, \text{BF16}\} \times \{\text{Buf 1}, \text{Buf 2}\} \times \{\text{False}, \text{semi}, \text{full}\}$) passed with zero staging leaks and bit-exact resumed trajectories (\texttt{atol} $= 0.0$).
\end{itemize}
\end{ReferenceCase}

\section{Stage III: Short-Run Proxy Model Screening}

Stage III tests candidate configurations on smaller proxy neural networks before committing to full training runs. By running short training schedules on proxy models, teams filter out unstable settings and verify low-precision training dynamics early.

\subsection{Proxy Model Selection \& Step Budget}
The engineering team selects a smaller proxy model representative of the target architecture (e.g., a multi-layer causal Transformer with Rotary Position Embeddings (RoPE), SwiGLU MLP, and RMSNorm):
\begin{enumerate}[label=\textbf{M\arabic*.}, itemsep=3pt]
    \item \textbf{Step Limit:} The screening run is capped at $5\%$ to $15\%$ of the full training budget (typically $500$ to $2{,}000$ optimization steps).
    \item \textbf{Batch Size Diversity:} Evaluate at least two effective batch sizes (standard $1\times$ micro-batch and $4\times$ gradient-accumulated batch) to check update stability under different gradient noise levels.
    \item \textbf{Controlled Schedule:} Use an identical linear warmup ($5\%$ of steps) followed by cosine decay to $5\%$ of peak learning rate across all candidates.
\end{enumerate}

\subsection{Early Stopping \& Rejection Criteria}
A candidate run is stopped and rejected immediately if any of the following three failure modes occur:
\begin{enumerate}[label=\textbf{F\arabic*.}, itemsep=3pt]
    \item \textbf{Update-to-Weight Ratio Breach:} Compute the layer-wise relative update norm across weight matrices ($\mathrm{ndim} \ge 2$, such as attention projections \texttt{q\_proj} and MLP gates \texttt{gate\_proj}, excluding 1D biases and normalization scales):
    \begin{equation}
        R_{t, l} = \frac{\|\delta_{t, l}\|_2}{\|\theta_{t, l}\|_2}.
    \end{equation}
    \textbf{Rejection Rule:} If $R_{t, l} > 0.10$ for any matrix layer $l$ after warmup (excessive step size) or $R_{t, l} < 10^{-7}$ during active training ($t > t_{\text{warmup}}$ and $\eta_t \ge 0.1 \eta_{\max}$, indicating stalled updates), the configuration is rejected.
    \item \textbf{Unrecovered Loss Spikes:} If the exponential moving average (EMA, $\alpha = 0.99$) training loss jumps by more than $0.5$ nats above its pre-spike minimum and fails to recover within $50$ consecutive steps, the run is rejected.
    \item \textbf{Low-Precision Stalling (Sub-ULP Update Cancellation):} Compare the final EMA training loss of the native low-precision (\texttt{bfloat16} or \texttt{float16}) proxy trial against an identical reference trial executed in full \texttt{float32} precision~\cite{gupta2015limited, zamirai2021bfloat16}:
    \begin{equation}
        \Delta\mathcal{L}_{\text{stall}} = \mathcal{L}_{\text{EMA}}^{(\text{BF16/FP16})}(T) - \mathcal{L}_{\text{EMA}}^{(\text{FP32-Ref})}(T).
    \end{equation}
    \textbf{Rejection Rule:} If $\Delta\mathcal{L}_{\text{stall}} \ge 0.25$ nats, the update rule suffers from low-precision mantissa truncation and is rejected.
\end{enumerate}

\subsection{Learning Rate \& Bounding Mode Grid}
Because bounded update velocities interact with the base learning rate, engineers evaluate a learning rate grid ($\eta \in [8\times 10^{-4}, 4\times 10^{-2}]$) paired with bounding modes (\texttt{bound\_\allowbreak type="coordinate"} vs.\ \texttt{"radial"}, default \texttt{bound\_\allowbreak ratio=0.05}):
\begin{itemize}[noitemsep,topsep=2pt]
    \item \textbf{Long-Run Setting ($>20\text{k}$ steps):} Identify conservative learning rates ($\eta \in [3\times 10^{-4}, 8\times 10^{-4}]$, $\lambda_{\text{wd}} = 8\times 10^{-2}$) where single-buffer and dual-buffer modes maintain stable update ratios.
    \item \textbf{Short-Run Setting ($<10\text{k}$ steps):} Identify higher learning rates ($\eta = 10^{-2}$) kept stable by coordinate-wise hyperbolic bounding~\cite{brock2021agc}.
\end{itemize}

\begin{ReferenceCase}{LuminaV (v1.5.3--v1.5.5) Proxy Model Screening in Stage III}
Stage III screening was executed on \texttt{MProxyLLM} (4-layer causal Transformer with RoPE, SwiGLU, and RMSNorm; $V = 2{,}048, d_{\text{model}} = 256, n_{\text{head}} = 4, d_{\text{ffn}} = 682$; $500$ steps, $25$-step warmup, batch size $8$, sequence length $128$) across $1\times$ and $4\times$ gradient accumulation:

\begin{center}
\resizebox{\linewidth}{!}{%
\begin{tabular}{@{}lccccccc@{}}
\toprule
\textbf{Configuration} & \textbf{BF16 EMA} & \textbf{FP32 Ref} & \textbf{Max $R_{t,l}$} & \textbf{Min $R_{t,l}$} & $\mathbf{F_1}$ & $\mathbf{F_2}$ & $\mathbf{F_3}$ \textbf{(Verdict)} \\
\midrule
\textbf{LuminaV Buf 1 ($1\times$, LR $8\times 10^{-4}$)} & $0.9018$ & $0.9021$ & $0.01574$ & $1.24\times 10^{-3}$ & PASS & PASS & PASS (\textbf{QUALIFIED}) \\
\textbf{LuminaV Buf 1 ($4\times$, LR $8\times 10^{-4}$)} & $\mathbf{0.8462}$ & $\mathbf{0.8462}$ & $0.01602$ & $1.24\times 10^{-3}$ & PASS & PASS & PASS (\textbf{QUALIFIED}) \\
\textbf{LuminaV Buf 2 ($1\times$, LR $8\times 10^{-4}$)} & $0.9146$ & $0.9141$ & $0.01604$ & $1.03\times 10^{-3}$ & PASS & PASS & PASS (\textbf{QUALIFIED}) \\
\textbf{LuminaV Buf 2 ($4\times$, LR $8\times 10^{-4}$)} & $0.8518$ & $0.8515$ & $0.01684$ & $9.81\times 10^{-4}$ & PASS & PASS & PASS (\textbf{QUALIFIED}) \\
\textbf{LuminaV Buf 1 ($1\times$, LR $1\times 10^{-2}$)} & $0.9326$ & $0.9317$ & $0.04041$ & $3.03\times 10^{-3}$ & PASS & PASS & PASS (\textbf{QUALIFIED}) \\
AdamW Native BF16 ($1\times$, LR $8\times 10^{-4}$) & $0.9763$ & $0.9552$ & $0.01423$ & $7.76\times 10^{-5}$ & PASS & PASS & PASS (QUALIFIED) \\
AdamW Native BF16 ($4\times$, LR $8\times 10^{-4}$) & $0.8767$ & $0.8678$ & $0.01600$ & $7.50\times 10^{-5}$ & PASS & PASS & PASS (QUALIFIED) \\
AdamW FP32 Master ($1\times$, LR $3\times 10^{-4}$) & $1.0989$ & $1.0992$ & $0.00611$ & $1.00\times 10^{-4}$ & PASS & PASS & PASS (QUALIFIED) \\
AdamW FP32 Master ($4\times$, LR $3\times 10^{-4}$) & $0.9693$ & $0.9692$ & $0.00708$ & $1.00\times 10^{-4}$ & PASS & PASS & PASS (QUALIFIED) \\
\bottomrule
\end{tabular}%
}
\end{center}

\textbf{Key Findings:} All five LuminaV configurations achieved $|\mathcal{L}_{\text{BF16}} - \mathcal{L}_{\text{FP32}}| \le 0.0009$ nats (zero BF16 precision drift without master weights), compared to a $+0.0211$ nat BF16 penalty on Native BF16 AdamW. Even at an aggressive learning rate of $\eta = 10^{-2}$, coordinate bounding (\texttt{bound\_ratio=0.05}) kept $\max R_{t,l} = 0.04041 \le 0.10$.
\end{ReferenceCase}

\section{Stage IV: Large-Scale Pre-Training Evaluation}

Stage IV is the final empirical validation phase, divided into two testing tiers:
\begin{enumerate}[noitemsep,topsep=2pt]
    \item \textbf{Stage IV-A (Hardware Qualification, Auto-Router \& Burn-in Check):} Medium-length runs ($2{,}500$ to $5{,}000$ steps, context length $512$) with gradient accumulation ($8\times$, $4{,}096$ tokens/step) measuring active \& isolated VRAM usage, kernel latency, hardware auto-router efficiency (\texttt{GPU\_\allowbreak DISPATCH\_\allowbreak THRESHOLD} $= 16{,}384$ elements), FSDP collective sharding parity, and held-out validation perplexity.
    \item \textbf{Stage IV-B (Full-Schedule \& Long-Context Qualification):} Full training runs ($\ge 20{,}000$ steps, context length $\ge 2{,}048$) checking long-horizon learning rate decay stability, gradient spike resilience, and held-out validation perplexity across the full schedule.
\end{enumerate}

\subsection{Data Pipeline Hygiene \& Cryptographic Zero-Leakage Protocol}
To prevent train--test overlap when streaming from massive web and bilingual corpora, LSO-Eval Stage IV enforces three data hygiene invariants prior to training:
\begin{enumerate}[label=\textbf{D\arabic*.}, noitemsep, topsep=2pt]
    \item \textbf{Cryptographic Hash Locking (XXH3-64 WAL):} All documents used to construct the held-out validation set ($N_{\text{eval}} \ge 450{,}000$ tokens) must be hashed via 64-bit XXH3 and committed to a persistent Write-Ahead Log (\texttt{.dedup} WAL).
    \item \textbf{Cross-Context Multiset Identity:} When evaluating across multiple sequence lengths (e.g., $880 \times 512 = 450{,}560$ tokens in Stage IV-A and $220 \times 2{,}048 = 450{,}560$ tokens in Stage IV-B), the flattened sorted token multisets must be $100.00\%$ bit-identical.
    \item \textbf{Strict Stream Offset \& WAL Inheritance:} Training prefetchers must skip all raw stream records consumed during validation construction (\texttt{skip\_raw\_items}) and preload the validation WAL into the live \texttt{StreamDeduplicator} so $100.00\%$ of validation hashes remain blocked throughout Stage IV-A and IV-B training.
\end{enumerate}

\subsection{Hardware Profiling, Isolated Optimizer Footprint \& Distributed Sharding}
All candidate and baseline optimizers run on identical hardware nodes with matched random seeds, data loaders, and sequence lengths. Under PyTorch Fully Sharded Data Parallel (FSDP) with \texttt{fsdp\_\allowbreak sync=True}, local reduction values $\big[\Sigma m_{\text{mask}}, \Sigma u_{\text{eff}}^2, \Sigma \theta_0^2, N_{\text{local}}\big]$ are packed into a single asynchronous \texttt{all\_\allowbreak reduce} call per parameter group.

Following memory-efficient optimizer literature (8-bit Optimizers~\cite{dettmers2022eightbit} and GaLore~\cite{zhao2024galore}), Stage IV decomposes peak active PyTorch tensor memory (\texttt{torch.cuda.\allowbreak max\_\allowbreak memory\_\allowbreak allocated()}) into the fixed non-optimizer baseline $\text{VRAM}_{\text{base}}$ and the \textbf{Isolated Optimizer Footprint} $\text{VRAM}_{\text{opt}}$ (Table~\ref{tab:stage4_dimensions}):
\begin{equation}
    \text{VRAM}_{\text{peak}} = \underbrace{W_{\text{model}} + G_{\text{grad}} + A_{\text{ctx}}}_{\text{Fixed PyTorch Baseline } (\text{VRAM}_{\text{base}})} + \underbrace{\text{VRAM}_{\text{opt}}}_{\text{Isolated Optimizer State \& Master Copy}}.
    \label{eq:vram_decomposition}
\end{equation}
Note that \texttt{max\_\allowbreak memory\_\allowbreak allocated()} measures active tensor bytes managed by PyTorch's caching allocator; it excludes unoccupied cached blocks (\texttt{max\_\allowbreak memory\_\allowbreak reserved()}) as well as CUDA driver contexts, cuBLAS workspaces, and Triton module binaries ($\approx 1.5$--$3.5\text{ GB}$ on \texttt{nvidia-smi}). Because $\text{VRAM}_{\text{base}}$ (model weights $W_{\text{model}}$, gradients $G_{\text{grad}}$, and autograd activations $A_{\text{ctx}}$) is fixed by the model architecture and sequence length, isolating $\text{VRAM}_{\text{opt}} = \text{VRAM}_{\text{peak}} - \text{VRAM}_{\text{base}}$ measures the true memory efficiency of the optimizer itself~\cite{dettmers2022eightbit, zhao2024galore}.

\begin{table}[!htbp]
\centering
\caption{Stage IV: Standard Evaluation Metrics}
\label{tab:stage4_dimensions}
\vspace{0.6em}
\renewcommand{\arraystretch}{1.2}
\begin{tabularx}{\textwidth}{@{}l l >{\raggedright\arraybackslash}X@{}}
\toprule
\textbf{Dimension} & \textbf{Metric} & \textbf{Measurement Method \& Tooling} \\
\midrule
\textbf{1. Convergence} & Raw Loss \& EMA Loss & Recorded at every step; smoothed via EMA ($\alpha \in \{0.99, 0.995\}$). \\
\textbf{2. Generalization} & Held-Out Validation PPL & Evaluated on $\ge 450\text{k}$ zero-leakage validation tokens ($\exp(\mathcal{L}_{\text{val}})$). \\
\textbf{3. Total \& Opt Memory} & $\text{VRAM}_{\text{peak}}$ \& $\text{VRAM}_{\text{opt}}$ (GB) & Total active tensor peak via \texttt{max\_memory\_allocated()}~\cite{paszke2019pytorch} and isolated $\text{VRAM}_{\text{opt}} = \text{VRAM}_{\text{peak}} - \text{VRAM}_{\text{base}}$~\cite{dettmers2022eightbit, zhao2024galore}. \\
\textbf{4. Latency} & Step \& Opt Kernel Time (ms) & End-to-end step time and isolated CUDA Event timer around \texttt{opt.step()}. \\
\bottomrule
\end{tabularx}
\end{table}

\subsection{Two-Tiered Pareto Acceptance \& Held-Out Perplexity Criteria}
Stage IV evaluates checkpoints on the held-out validation set:
\begin{equation}
    \text{PPL}_{\text{held-out}} = \exp\left( \frac{1}{N} \sum_{i=1}^N \mathcal{L}_{\text{CE}}(x_i) \right).
\end{equation}
Because single-buffer ($2\text{ B/param}$) and dual-buffer ($4\text{--}12\text{ B/param}$) optimizers target distinct hardware operating regimes~\cite{dettmers2022eightbit, zhao2024galore}, LSO-Eval v1.1.0 enforces a \textbf{Two-Tiered Pareto Acceptance Rule} relative to the reference Native BF16 AdamW baseline~\cite{loshchilov2019adamw}:
\begin{enumerate}[label=\textbf{Class \Roman*.}, itemsep=3pt]
    \item \textbf{Dual-Buffer / Full-State Optimizers ($\ge 4\text{ bytes/param}$, Quality-Dominant Tier):} The candidate maintains both first-moment and second-moment buffers ($m_t, v_t$). It must achieve \textbf{strictly equal or lower held-out validation perplexity} than the reference AdamW baseline at matched token count ($\Delta\text{PPL}_{\text{held-out}} = \text{PPL}_{\text{cand}} - \text{PPL}_{\text{AdamW-BF16}} \le 0.0$) while passing $\text{F}_1 \land \text{F}_2$ stability gates with zero extra VRAM overhead ($\Delta\text{VRAM}_{\text{opt}} \le 0.0\text{ GB}$).
    \item \textbf{Single-Buffer / State-Compressed Optimizers ($\le 2\text{ bytes/param}$, Memory-Pareto Tier):} The candidate eliminates the second-moment variance buffer $v_t$ entirely to operate on memory-constrained accelerators. It qualifies as \textbf{PASSED (Memory-Pareto Qualified)} if and only if it simultaneously satisfies all four conditions:
    \begin{enumerate}[label=(\alph*), noitemsep, topsep=1pt]
        \item \textit{Dynamic Stability ($\text{F}_1 \land \text{F}_2$):} Layer-wise update ratios remain strictly within $R_{t,l} \in [10^{-7}, 0.10]$ with zero unrecovered loss spikes across the full schedule.
        \item \textit{Monotonic Generalization:} Held-out validation perplexity $\text{PPL}_{\text{held-out}}(t)$ decreases monotonically across evaluation checkpoints without late-stage divergence or rebound.
        \item \textit{Hardware Footprint Dominance:} Reduces isolated optimizer state memory by at least $50\%$ ($\Delta\text{VRAM}_{\text{opt}} \le -50.0\%$, or total $\Delta\text{VRAM}_{\text{peak}} \le -15.0\%$) relative to Native BF16 AdamW (and $\Delta\text{VRAM}_{\text{opt}} \le -80.0\%$ relative to FP32 Master AdamW) while achieving faster per-step kernel latency ($T_{\text{opt}} < T_{\text{AdamW-BF16}}$).
        \item \textit{Bounded Compression Loss Tolerance:} The validation log-perplexity gap induced by omitting the second-moment buffer remains strictly within the Stage III $\text{F}_3$ tolerance bound:
        \begin{equation}
            \Delta\mathcal{L}_{\text{val}} = \ln\left(\frac{\text{PPL}_{\text{cand}}}{\text{PPL}_{\text{AdamW-BF16}}}\right) \le 0.25\text{ nats}.
        \end{equation}
    \end{enumerate}
\end{enumerate}

\begin{ReferenceCase}{LuminaV (v1.5.3--v1.5.5) 4B-Parameter Hybrid LLM Qualification in Stage IV-A \& IV-B}
\textbf{Version Invariance Note:} The Stage IV-A and IV-B empirical qualification benchmarks below were executed using the \texttt{LuminaV v1.5.3} release binary in uncentered defect mode (\texttt{centered\_defect=False}, $d_t = g_t - m$) on native \texttt{torch.bfloat16} parameters (\texttt{dtype\_mode=1}). Because \texttt{v1.5.4} modified only pure \texttt{float16} (\texttt{dtype\_mode=2}) state containers and removed the unused centered-defect branch, and \texttt{v1.5.5} added opt-in contiguous bucket coalescing (\texttt{fuse\_tensors=False} by default) and initialization-time \texttt{lr=0.0} JIT warmup, the \texttt{v1.5.3} native \texttt{bfloat16} update equations and kernel instructions are \textbf{100\% mathematically and bit-for-bit trajectory-identical} to \texttt{v1.5.4} and \texttt{v1.5.5}.

\smallskip
\noindent\textbf{1. Model Architecture \& Cryptographic Zero-Leakage Proof:}
Evaluations were conducted from scratch ($100\%$ random initialization) on the hybrid \textbf{Qwen/Qwen3.5-4B-Base} architecture ($32$ layers: $24$ GatedDeltaNet linear-attention layers with \texttt{causal\_conv1d} + $8$ full SDPA layers; $\approx 4.26\times 10^9$ parameters, fixed model weight + gradient footprint $W_{\text{BF16}} + G_{\text{BF16}} = 15.88\text{ GB}$). The data pipeline streamed from a 6-way interleaved English--Indonesian corpus filtered via Rust \texttt{FusedPipeline}, \texttt{PIIMasker}, and \texttt{StreamDeduplicator}:
\begin{itemize}[noitemsep,topsep=1pt]
    \item \textit{Validation Tensor Parity:} Stage IV-A ($880 \times 512 = 450{,}560$ tokens, SHA-256 \texttt{4d278cb8\allowbreak 7e4abbe1}) and Stage IV-B ($220 \times 2{,}048 = 450{,}560$ tokens, SHA-256 \texttt{108e517e\allowbreak 1c181cb6}) consumed exactly $622$ raw stream records (\texttt{raw\_\allowbreak offset}) and verified \textbf{100.00\% identical token multisets}.
    \item \textit{Zero-Leakage WAL Continuity:} All $143{,}475$ XXH3-64 validation hashes from Stage IV-A remained $100.00\%$ locked inside the $295{,}605$ committed hashes of the Stage IV-B WAL (\texttt{PASSED: 100\% Bit-Exact Eval \& Zero Leakage}).
\end{itemize}

\smallskip
\noindent\textbf{2. Stage IV-A Qualification (Context $512$, Grad Accum $8\times$, $4{,}800$ Steps / $19.66\text{M}$ Tokens):}
At sequence length $512$, the fixed PyTorch non-optimizer baseline (weights + gradients + autograd activations) is $\text{VRAM}_{\text{base}}^{(512)} \approx 20.87\text{ GB}$. Across the $4{,}800$-step accumulated qualification run evaluated on $450{,}560$ held-out tokens:

\begin{center}
\resizebox{\linewidth}{!}{%
\begin{tabular}{@{}lccccccccc@{}}
\toprule
\textbf{Optimizer Configuration} & \textbf{Final EMA} & \textbf{Val PPL} & $\bm{\Delta}$\textbf{PPL} & \textbf{Peak VRAM} & \textbf{Opt VRAM ($\text{VRAM}_{\text{opt}}$)} & \textbf{Max $R_{t,l}$} & \textbf{Step Time} & \textbf{Opt Time} & \textbf{LSO-Eval Status} \\
\midrule
\textbf{LuminaV Buffer 2 (Dual, $8\times 10^{-4}$)} & $\mathbf{4.7621}$ & $\mathbf{141.96}$ & $\mathbf{-39.21}$ & $36.75\text{ GB}$ & $15.88\text{ GB}$ ($0.0\%$) & $0.0047$ & $841.6\text{ ms}$ & $48.37\text{ ms}$ & \textbf{PASSED (Class I)} \\
\textbf{LuminaV Buffer 1 (TD-CLM, $3\times 10^{-4}$)} & $5.2569$ & $195.58$ & $+14.41$ & $\mathbf{28.81\text{ GB}}$ & $\mathbf{7.94\text{ GB}}$ ($\mathbf{-50.0\%}$) & $0.0120$ & $\mathbf{829.4\text{ ms}}$ & $36.77\text{ ms}$ & \textbf{PASSED (Class II)} \\
\textbf{LuminaV Buffer 1 (TD-CLM, $8\times 10^{-4}$)} & $5.2953$ & $221.07$ & $+39.90$ & $\mathbf{28.81\text{ GB}}$ & $\mathbf{7.94\text{ GB}}$ ($\mathbf{-50.0\%}$) & $0.0191$ & $830.7\text{ ms}$ & $\mathbf{36.46\text{ ms}}$ & \textbf{PASSED (Class II)} \\
AdamW (Native BF16, $3\times 10^{-4}$)~\cite{loshchilov2019adamw} & $5.4497$ & $181.17$ & $+0.00$ & $36.75\text{ GB}$ & $15.88\text{ GB}$ (Ref) & $0.0034$ & $835.4\text{ ms}$ & $41.58\text{ ms}$ & Baseline (BF16) \\
AdamW (FP32 Master, $3\times 10^{-4}$)~\cite{loshchilov2019adamw} & $5.0234$ & $153.78$ & $-27.39$ & $75.36\text{ GB}$ & $54.49\text{ GB}$ ($+243\%$) & $0.0040$ & $919.5\text{ ms}$ & $123.35\text{ ms}$ & Baseline (Master) \\
\bottomrule
\end{tabular}%
}
\end{center}

\smallskip
\noindent\textbf{3. Stage IV-B Full-Schedule Qualification (Context $2{,}048$, Grad Accum $1\times$, $20{,}000$ Steps / $40.96\text{M}$ Tokens):}
At sequence length $2{,}048$, autograd activations expand the fixed non-optimizer baseline to $\text{VRAM}_{\text{base}}^{(2048)} \approx 28.03\text{ GB}$, while isolated optimizer state allocations remain invariant at $\mathbf{7.94\text{ GB}}$ (Buffer 1) and $\mathbf{15.86\text{ GB}}$ (Buffer 2):

\begin{center}
\resizebox{\linewidth}{!}{%
\begin{tabular}{@{}lccccccccc@{}}
\toprule
\textbf{Optimizer Configuration} & \textbf{Final EMA} & \textbf{Val PPL} & $\bm{\Delta}$\textbf{PPL} & \textbf{Peak VRAM} & \textbf{Opt VRAM ($\text{VRAM}_{\text{opt}}$)} & \textbf{Max $R_{t,l}$} & \textbf{Step Time} & \textbf{Opt Time} & \textbf{LSO-Eval Status} \\
\midrule
\textbf{LuminaV Buffer 2 (Dual, $8\times 10^{-4}$)} & $\mathbf{4.5549}$ & $\mathbf{89.14}$ & $\mathbf{-43.43}$ & $43.89\text{ GB}$ & $15.86\text{ GB}$ ($0.0\%$) & $0.0189$ & $297.1\text{ ms}$ & $48.31\text{ ms}$ & \textbf{PASSED (Class I)} \\
\textbf{LuminaV Buffer 1 (TD-CLM, $3\times 10^{-4}$)} & $4.9173$ & $165.13$ & $+32.56$ & $\mathbf{35.97\text{ GB}}$ & $\mathbf{7.94\text{ GB}}$ ($\mathbf{-50.0\%}$) & $0.0087$ & $\mathbf{283.5\text{ ms}}$ & $\mathbf{36.99\text{ ms}}$ & \textbf{PASSED (Class II)} \\
AdamW (Native BF16, $3\times 10^{-4}$)~\cite{loshchilov2019adamw} & $4.7621$ & $132.57$ & $+0.00$ & $43.89\text{ GB}$ & $15.86\text{ GB}$ (Ref) & $0.0037$ & $288.0\text{ ms}$ & $41.48\text{ ms}$ & Baseline (BF16) \\
AdamW (FP32 Master, $3\times 10^{-4}$)~\cite{loshchilov2019adamw} & $4.8116$ & $124.43$ & $-8.14$ & $80.11\text{ GB}$ & $52.08\text{ GB}$ ($+228\%$) & $0.0045$ & $370.1\text{ ms}$ & $123.33\text{ ms}$ & Baseline (Master) \\
\bottomrule
\end{tabular}%
}
\end{center}

\textbf{Key Qualification Findings:}
\begin{enumerate}[noitemsep,topsep=1pt]
    \item \textbf{Class I Quality Dominance (LuminaV Buffer 2):} Without using FP32 master weights, \texttt{LuminaV Buffer 2 (Dual)} achieved the lowest validation perplexity in both Stage IV-A ($\mathbf{141.96}$, beating Native BF16 AdamW by $-39.21$ PPL and FP32 Master AdamW by $-11.82$ PPL) and Stage IV-B ($\mathbf{89.14}$, beating Native BF16 AdamW by $\mathbf{-43.43}$ \textbf{PPL} and outperforming FP32 Master AdamW's $124.43$ by $\mathbf{-35.29}$ \textbf{PPL} while saving $\mathbf{36.22\text{ GB}}$ of VRAM and running $2.55\times$ faster in optimizer step latency).
    \item \textbf{Class II Memory-Pareto Qualification (LuminaV Buffer 1):} \texttt{LuminaV Buffer 1 (TD-CLM, LR $3\times 10^{-4}$)} reduced isolated optimizer state memory by exactly $\mathbf{50.0\%}$ ($7.94\text{ GB}$ vs.\ $15.88\text{ GB}$ on BF16 AdamW) and reduced optimizer memory overhead by $\mathbf{84.8\%\text{--}85.4\%}$ ($7.94\text{ GB}$ vs.\ $52.08\text{--}54.49\text{ GB}$ on FP32 Master AdamW, saving $44.14\text{--}46.55\text{ GB}$ of total HBM so 4B pre-training fits comfortably on a single $40\text{ GB}$ GPU where FP32 Master AdamW fails with OOM). Its kernel latency ($36.46\text{--}36.99\text{ ms}$) was $3.33\times$ faster than FP32 Master AdamW ($123.33\text{ ms}$), while its validation log-perplexity gap remained strictly bounded within Class II tolerances ($\Delta\mathcal{L}_{\text{val}} = 0.0765 \le 0.25$ nats in IV-A and $0.2196 \le 0.25$ nats in IV-B).
\end{enumerate}
\end{ReferenceCase}

\section{Differential Versioning \& Regression Protocol (D-Eval)}

When updating an existing optimizer or backend kernel from version $v_{k-1}$ to $v_k$, LSO-Eval requires the \textbf{D-Eval Protocol} to ensure no silent performance or numerical regressions are introduced.

\subsection{The Four Anti-Regression Gates}
A release candidate must pass all four gates before deployment (Table~\ref{tab:deval_gates}):

\begin{table}[!htbp]
\centering
\caption{The D-Eval Four-Gate Anti-Regression Framework}
\label{tab:deval_gates}
\vspace{0.4em}
\renewcommand{\arraystretch}{1.25}
\begin{tabularx}{\textwidth}{@{}l l >{\raggedright\arraybackslash}X l@{}}
\toprule
\textbf{Gate} & \textbf{Evaluation Target} & \textbf{Acceptance Criterion} & \textbf{Verification Tool} \\
\midrule
\textbf{Gate 1} & Algebraic Safety & All 12 Stage I SMT queries return \texttt{UNSAT}; Pad\'{e} \& SR bounds preserved. & Z3~\cite{demoura2008z3} / SymPy~\cite{meurer2017sympy} \\
\textbf{Gate 2} & Trajectory \& Backend Parity & Head-to-head 2D run on $\kappa \ge 1000$ and cross-backend parity tests pass. & 2D Sim \& Rust Tests \\
\textbf{Gate 3} & Checkpoint Compatibility & Checkpoints saved in version $v_{k-1}$ load cleanly into $v_k$ with identical dtypes and trajectories. & State-Dict Audit \\
\textbf{Gate 4} & Production Metrics Delta & Meets non-regression bounds: $\Delta \mathcal{L} \le 0$, $\Delta \text{VRAM}_{\text{opt}} \le 0$, $\Delta T_{\text{kernel}} \le 0$. & 4B Scale Benchmark \\
\bottomrule
\end{tabularx}
\end{table}

\subsection{Performance Delta Comparison}
For Gate 4, the engineering team reports the performance differences against the previous version:
\begin{align}
    \Delta \mathcal{L}_{\text{EMA}} &= \mathcal{L}_{\text{EMA}}(v_k) - \mathcal{L}_{\text{EMA}}(v_{k-1}) \le 0, \\
    \Delta \text{VRAM}_{\text{opt}} &= \text{VRAM}_{\text{opt}}(v_k) - \text{VRAM}_{\text{opt}}(v_{k-1}) \le 0, \\
    \Delta T_{\text{kernel}} &= T_{\text{kernel}}(v_k) - T_{\text{kernel}}(v_{k-1}) \le 0.
\end{align}
If any metric regresses ($\Delta > 0$), the team must document the trade-off in an Architecture Decision Record (ADR) explaining why the change is justified.

\begin{ReferenceCase}{D-Eval Audit: LuminaV Evolution from v1.0.0 to v1.5.3--v1.5.5}
Summary of the D-Eval regression audits across the LuminaV release series (\texttt{v1.0.0} $\to$ \texttt{v1.5.3} $\to$ \texttt{v1.5.5}):
\begin{itemize}[itemsep=2pt]
    \item \textbf{Gate 1 (SMT \& Unified $\epsilon_{\text{eff}}$):} Unified Convex Lookahead Momentum weights ($w_1 = \beta_1^2, w_2 = 1-\beta_1^2$) across Buffer 1 and Buffer 2 ($w_1 + w_2 = 1$), verified $[13/6]$ Pad\'{e} $\tanh$ accuracy ($9.48 \times 10^{-8} < 1.5 \times 10^{-6}$), and in \texttt{v1.5.4} promoted pure \texttt{FP16} 16-bit state containers to in-container Bitwise-BF16 ($8$-bit exponent range $1.18\times 10^{-38}$ at zero extra VRAM), unifying $\epsilon_{\text{eff}} = \epsilon = 10^{-8}$ across all precisions.
    \item \textbf{Gate 2 (Uncentered Defect Simplification \& Backend Parity):} Removed the global defect mean reduction ($\mu_d$) in \texttt{v1.5.4} to standardize on the local coordinate-wise innovation defect $d_t^{(\text{B1})} = g_t - m_{t-1}$ and $d_t^{(\text{B2})} = g_t - m_t$, eliminating pre-step reduction kernels while preserving exact trajectory parity with \texttt{v1.5.3} (\texttt{centered\_defect=False}) across Scalar, AVX2, AVX-512, Vulkan 1.3/1.4 Wave32 UMA, and Triton AOT/JIT.
    \item \textbf{Gate 3 (Contiguous Bucket Fusion \& Checkpoint Cleanliness):} Added opt-in bounded contiguous parameter coalescing (\texttt{fuse\_\allowbreak tensors=False}, \texttt{bucket\_\allowbreak cap\_\allowbreak mb=500.0}) and zero-stall Triton JIT pre-warmup (\texttt{prewarm\_\allowbreak jit=True}) in \texttt{v1.5.5}, while ensuring the 6-slot \texttt{float32} scratchpad accumulators ($s_0, s_1, s_2, s_5$) are zeroed on export without leaking staging keys.
    \item \textbf{Gate 4 (Zero-Overhead Hardware Footprint):} Maintained $\Delta \text{VRAM}_{\text{opt}} = 0\text{ GB}$ ($7.94\text{ GB}$ isolated optimizer state on 4B LLM for Buffer 1; $15.86\text{ GB}$ for Buffer 2) and packaged zero-dependency sub-1MB CPython Stable ABI (\texttt{abi3-py310} / \texttt{abi3t}) wheels.
\end{itemize}
\end{ReferenceCase}

\section{Standard LSO-Eval Certification Checklist}

No optimizer or hardware kernel update may be merged into main production branches until every item in Table~\ref{tab:certification_checklist} is verified and signed off.

\begin{table}[!htbp]
\centering
\caption{Official LSO-Eval v1.1.0 Release Certification Checklist}
\label{tab:certification_checklist}
\vspace{0.3em}
\small
\renewcommand{\arraystretch}{1.15}
\begin{tabularx}{\textwidth}{@{}c l >{\raggedright\arraybackslash}p{4.1cm} >{\raggedright\arraybackslash}X c@{}}
\toprule
\textbf{Stage} & \textbf{ID} & \textbf{Verification Target} & \textbf{Acceptance Criterion} & \textbf{Sign-off} \\
\midrule
\multirow{4}{*}{\textbf{Stage I}}
 & I.1 & 12-Query SMT Safety Suite & All 12 negated Z3 safety queries return \texttt{UNSAT}. & $[\quad]$ \\
 & I.2 & Large-Gradient Limit & Profile A: $\lim_{|g|\to\infty} |u_t| < 1.0$; Profile B/C: $\lim_{|g|\to\infty} \|u_t\| < \infty$. & $[\quad]$ \\
 & I.3 & Pad\'{e} $\tanh$ \& Continuity & $\|\tanh_{\text{Pad\'{e}}} - \tanh\|_\infty < 1.5\times 10^{-6}$; radial jump $\Delta_{\text{jump}} < \epsilon_{\text{FP32}}$. & $[\quad]$ \\
 & I.4 & Stochastic Rounding & Unbiased rounding ($|\mu_{\mathcal{Q}} - x| < 10^{-6}$, $p_{\text{TOST}} < 0.01$ or $t$-test $p \ge 0.01$) over $200\text{k}$ samples. & $[\quad]$ \\
\midrule
\multirow{5}{*}{\textbf{Stage II}}
 & II.1 & Synthetic Test Surfaces & Stable convergence on $\kappa \ge 1000$ (B1), Rosenbrock (B2), Noise (B3), Rastrigin (B4), and $N=100$ Saddle (B5). & $[\quad]$ \\
 & II.2 & Balance Calibration & Default $\lambda = 0.30$ ($\lambda^* = 0.3008$) and $\pm 30\%$ hyperparameter perturbations degrade steps by $\le 15\%$. & $[\quad]$ \\
 & II.3 & 1-Pass EMA Parity & Fused 1-Pass Scratchpad EMA matches 2-Pass steps within $+1.0\%$ ($\alpha_{\text{ema}} \in \{0.95, 0.80\}$). & $[\quad]$ \\
 & II.4 & Cross-Backend Parity & Scalar, AVX2, AVX-512, Vulkan Wave32 UMA, Triton, and NumPy/PyTorch parity verified. & $[\quad]$ \\
 & II.5 & Checkpoint Roundtrip & No staging keys; $s_{0,1,2,5}$ zeroed; FP32/BF16 state dtypes \& bit-exact trajectory preserved on reload. & $[\quad]$ \\
\midrule
\multirow{3}{*}{\textbf{Stage III}}
 & III.1 & Update Norm Stability ($\text{F}_1$) & Matrix ($\mathrm{ndim} \ge 2$) active-phase update ratio satisfies $R_{t, l} \in [10^{-7}, 0.10]$. & $[\quad]$ \\
 & III.2 & Loss Spike Recovery ($\text{F}_2$) & Zero unrecovered EMA loss spikes ($>0.5$ nats for $>50$ steps) over proxy schedule. & $[\quad]$ \\
 & III.3 & Low-Precision Stall Gate ($\text{F}_3$) & Native BF16/FP16 final EMA loss within $< 0.25$ nats of matched FP32 reference trial. & $[\quad]$ \\
\midrule
\multirow{4}{*}{\textbf{Stage IV}}
 & IV.0 & Zero-Leakage WAL Proof & $100\%$ token multiset parity across contexts and $100\%$ XXH3-64 validation hash locking in training WAL. & $[\quad]$ \\
 & IV-A & Burn-In \& Hardware Profile & Isolated $\text{VRAM}_{\text{opt}}$, total $\text{VRAM}_{\text{peak}}$, kernel latency, and $4{,}800$-step convergence verified. & $[\quad]$ \\
 & IV-B.1 & Full-Schedule Stability & Zero late-stage divergence or $\text{F}_1/\text{F}_2$ breaches across full schedule ($\ge 20\text{k}$ steps, context $\ge 2{,}048$). & $[\quad]$ \\
 & IV-B.2 & Two-Tiered Pareto PPL & Class I (Dual-Buffer): $\Delta\text{PPL}_{\text{held-out}} \le 0.0$; Class II (Single-Buffer): $\Delta\text{VRAM}_{\text{opt}} \le -50\%, \Delta\mathcal{L}_{\text{val}} \le 0.25$. & $[\quad]$ \\
\midrule
\textbf{D-Eval}
 & D.1 & 4-Gate Anti-Regression & Verified non-regression against predecessor version across Gates 1--4. & $[\quad]$ \\
\bottomrule
\end{tabularx}
\end{table}

\clearpage
\begin{thebibliography}{99}
\small
\sloppy
\setlength{\itemsep}{3pt}

\bibitem{demoura2008z3}
L. de Moura and N. Bj{\o}rner, ``Z3: An efficient SMT solver,'' in \textit{Proc. Tools and Algorithms for the Construction and Analysis of Systems (TACAS)}, Budapest, Hungary, 2008, pp. 337--340.

\bibitem{maligranda2008triangle}
L. Maligranda, ``Some remarks on the triangle inequality for norms,'' \textit{Banach J. Math. Anal.}, vol. 2, no. 2, pp. 31--41, 2008.

\bibitem{chen2023lion}
X. Chen \textit{et al.}, ``Symbolic discovery of optimization algorithms,'' in \textit{Adv. Neural Inf. Process. Syst. (NeurIPS)}, vol. 36, 2023, pp. 49205--49233.

\bibitem{brock2021agc}
A. Brock, S. De, S. L. Smith, and K. Simonyan, ``High-performance large-scale image recognition without normalization,'' in \textit{Proc. 38th Int. Conf. Mach. Learn. (ICML)}, vol. 139, 2021, pp. 1059--1071.

\bibitem{loshchilov2019adamw}
I. Loshchilov and F. Hutter, ``Decoupled weight decay regularization,'' in \textit{Proc. Int. Conf. Learn. Represent. (ICLR)}, New Orleans, LA, USA, 2019.

\bibitem{zhuang2020adabelief}
J. Zhuang, T. Tang, Y. Ding, S. Tatikonda, N. Dvornek, X. Papademetris, and J. S. Duncan, ``AdaBelief optimizer: Adapting stepsizes by the belief in observed gradients,'' in \textit{Adv. Neural Inf. Process. Syst. (NeurIPS)}, vol. 33, 2020, pp. 18795--18806.

\bibitem{taniguchi2024adopt}
S. Taniguchi, K. Harada, G. Minegishi, Y. Oshima, S. C. Jeong, G. Nagahara, T. Iiyama, M. Suzuki, Yusuke Iwasawa, and Y. Matsuo, ``ADOPT: Modified Adam can converge with any $\beta_2$ with the optimal rate,'' in \textit{Adv. Neural Inf. Process. Syst. (NeurIPS)}, vol. 37, 2024.

\bibitem{gupta2018shampoo}
V. Gupta, T. Koren, and Y. Singer, ``Shampoo: Preconditioned stochastic tensor optimization,'' in \textit{Proc. 35th Int. Conf. Mach. Learn. (ICML)}, vol. 80, Stockholm, Sweden, 2018, pp. 1842--1850.

\bibitem{zhao2024galore}
J. Zhao, Z. Zhang, B. Chen, Z. Wang, A. Anandkumar, and Y. Tian, ``GaLore: Memory-efficient LLM training by gradient low-rank projection,'' in \textit{Proc. 41st Int. Conf. Mach. Learn. (ICML)}, vol. 235, Vienna, Austria, 2024.

\bibitem{liang2026cautious}
K. Liang, L. Chen, B. Liu, and Q. Liu, ``Cautious optimizers: Improving training with one line of code,'' in \textit{Proc. Int. Conf. Learn. Represent. (ICLR)}, 2026.

\bibitem{zhang2019lookahead}
M. R. Zhang, J. Lucas, G. E. Hinton, and J. Ba, ``Lookahead optimizer: $k$ steps forward, 1 step back,'' in \textit{Adv. Neural Inf. Process. Syst. (NeurIPS)}, vol. 32, Vancouver, BC, Canada, 2019, pp. 9593--9604.

\bibitem{meurer2017sympy}
A. Meurer \textit{et al.}, ``SymPy: Symbolic computing in Python,'' \textit{PeerJ Comput. Sci.}, vol. 3, p. e103, Jan. 2017, doi: 10.1717/peerj-cs.103.

\bibitem{molina2020pade}
A. Molina, P. Schramowski, and K. Kersting, ``Pad\'{e} activation units: End-to-end learning of flexible activation functions in deep networks,'' in \textit{Proc. Int. Conf. Learn. Represent. (ICLR)}, 2020.

\bibitem{gupta2015limited}
S. Gupta, A. Agrawal, K. Gopalakrishnan, and P. Narayanan, ``Deep learning with limited numerical precision,'' in \textit{Proc. 32nd Int. Conf. Mach. Learn. (ICML)}, Lille, France, 2015, pp. 1737--1746.

\bibitem{zamirai2021bfloat16}
P. Zamirai, J. Zhang, C. R. Aberger, and C. De Sa, ``Revisiting BFloat16 training,'' \textit{arXiv preprint arXiv:2010.06192v2}, Mar. 2021.

\bibitem{jamil2013benchmark}
M. Jamil and X.-S. Yang, ``A literature survey of benchmark functions for global optimization problems,'' \textit{Int. J. Math. Model. Numer. Optim.}, vol. 4, no. 2, pp. 150--194, 2013, doi: 10.1504/IJMMNO.2013.055204.

\bibitem{rosenbrock1960}
H. H. Rosenbrock, ``An automatic method for finding the greatest or least value of a function,'' \textit{Comput. J.}, vol. 3, no. 3, pp. 175--184, Oct. 1960, doi: 10.1093/comjnl/3.3.175.

\bibitem{rastrigin1974}
L. A. Rastrigin, \textit{Systems of Extremal Control}. Moscow, Russia: Nauka, 1974.

\bibitem{dao2022flashattention}
T. Dao, D. Y. Fu, S. Ermon, A. Rudra, and C. R\'{e}, ``FlashAttention: Fast and memory-efficient exact attention with IO-awareness,'' in \textit{Adv. Neural Inf. Process. Syst. (NeurIPS)}, vol. 35, New Orleans, LA, USA, 2022, pp. 16344--16359.

\bibitem{tillet2019triton}
P. Tillet, H. T. Kung, and D. Cox, ``Triton: An intermediate language and compiler for tiled neural network computations,'' in \textit{Proc. 3rd ACM SIGPLAN Int. Workshop Mach. Learn. Program. Lang. (MAPL)}, Phoenix, AZ, USA, Jun. 2019, pp. 10--19, doi: 10.1145/3315508.3329973.

\bibitem{dettmers2022eightbit}
T. Dettmers, M. Lewis, S. Shleifer, and L. Zettlemoyer, ``8-bit optimizers via block-wise quantization,'' in \textit{Proc. Int. Conf. Learn. Represent. (ICLR)}, 2022.

\bibitem{paszke2019pytorch}
A. Paszke \textit{et al.}, ``PyTorch: An imperative style, high-performance deep learning library,'' in \textit{Adv. Neural Inf. Process. Syst. (NeurIPS)}, vol. 32, Vancouver, BC, Canada, 2019, pp. 8024--8035.

\end{thebibliography}

\vspace{1.5em}
\noindent\begin{minipage}{\textwidth}
\noindent\rule{\textwidth}{0.4pt}\\[0.4em]
\centering\small
\textbf{Lumina Moon Standard Optimizer Evaluation Methodology (LSO-Eval v1.1.0)}\\
\textit{Official Engineering Specification --- Lumina Moon (cloverx-id)}
\end{minipage}

\end{document}
